\chapter{Conclusion}
This rapport has explored the various uses for, and workings of, spectroscopy.
First off, the fact that hydrogen emits light of very specific wavelengths, 
has been strongly proven, as there has been an essentially perfect correspondance between the observed emission lines, and the Balmer series.
It was also possible to associate some of the dips in the solar spectrum, with emission lines from Hydrogen and Helium,
which suggests that these elements, in the suns atmosphere, are responsible for absorbing the lowered wavelengths.

The solar spectrum was also be used to estimate the temperature of sun as $T = 5564.1 \text{ K} \pm 12.6 \text{ K}$,
which is 3.6 percent off compared to the accepted value.

The absorbance spectra had significant noise, but identification of the different salts were still possible.
The accuracy of the relative concentrations is however uncertain, and we have no table values to compare to.

Finally, we built a spectrometer, which is considerably worse than the one from OceanOptics.
The spectrometer measured high intensities in the same 'regions' as the actual spectra, and goes to zero in the same 'regions' (except Hydrogen for some reason), 
but the location of the peaks are all offset by varying amounts, that makes the spectrometer useless as a measument device.